\subsection{Review of Finance}

\content{
        Formulas so far: \\
        \fbox{\begin{minipage}{\textwidth}
        For these formulas, we have the following: 
        $P_0$ is the principal. 
        $r$ is the interest rate. 
        $t$ is the number of time intervals (specified by problem) which
        have passed. 
        $A$ is the amount paid back.
        $I$ is the interest.

        \begin{center}
        \begin{tabular}{|c|c|}
        \hline
        \T\B Simple Interest &  $ A = P_0(1+r) $ \\
        \hline
        \T\B Simple Interest over Time & $  A = P_0(1+rt) $ \\
        \hline
        \end{tabular}
        \end{center}
        \end{minipage}}
        \fbox{\begin{minipage}{\textwidth}
        For these formulas, we have the following: 
        $P_0$ is the principal, or amount of a loan. 
        $r$ is the annual interest rate.
        $t$ is the number of years which have passed. 
        $A$ is the total amount.
        $k$ is the number of times interest is compounded per year. 
        $d$ is the amount we deposit (or withdraw) from the account regularly,
        or the payment amount on a loan.

        \begin{center}
        \begin{tabular}{|c|c|}
        \hline
        \T\B Compound Interest & $ A =
        P_0\left(1+\frac{r}{k}\right)^{kt} $ \\
        \hline
        \T\B Savings Annuity & $ A =
        d\frac{\left(\left(1+\frac{r}{k}\right)^{kt}-1\right)}{\frac{r}{k}} $\\
        \hline
        \T\B Payout Annuity/Loans & $P_0 =
        d\frac{\left(1-\left(1+\frac{r}{k}\right)^{-kt}\right)}{\frac{r}{k}}$ \\
        \hline
        \end{tabular}
        \end{center}
        \end{minipage}}
}{% presentation
}{% nopresentation
}{% comments
}

\content{
        \begin{example}
        Suppose you obtain a \$3,000 T-note with a 3\% annual rate, paid quarterly, with
        maturity in 5 years. How much interest will you earn?
        \end{example}
}{% presentation
\vspace{2in}
}{% nopresentation
\vspace{2in}
}{% comments
}

\content{
\begin{example}
You deposit \$300 in an account earning 5\% interest compounded annually. How
much will you have in the account in 10 years?
\end{example}
}{% presentation
\vspace{2in}
}{% nopresentation
\vspace{2in}
}{% comments
}

\content{
        \begin{example}
        You want to be able to withdraw \$30,000 each year for 25 years. Your account earns
        8\% interest.
        \begin{enumerate}
                \item How much do you need in your account at the beginning
                \item How much total money will you pull out of the account?
                \item How much of that money is interest?
        \end{enumerate}
        \end{example}
}{% presentation 
\vspace{3in}
}{% nopresentation
\vspace{3in}
}{% comments
}

\content{
        \begin{example}
        You want to buy a \$25,000 car. The company is offering a 2\% interest rate for 48
        months (4 years). What will your monthly payments be?
        \end{example}
}{% presentation
\vspace{3in}
}{% nopresentation
\vspace{3in}
}{% comments
}

\content{
        \begin{example}
        A friend lends you \$200 for a week, which you agree to repay with 5\% one-time
        interest. How much will you have to repay?
        \end{example}
}{% presentation
\vspace{2in}
}{% nopresentation
\vspace{2in}
}{% comments
}

\content{
        \begin{example}
        You wish to have \$3000 in 2 years to buy a fancy new stereo system. How much
        should you deposit each quarter into an account paying 8\% compounded quarterly?        
        \end{example}
}{% presentation
\vspace{3in}
}{% nopresentation
\vspace{3in}
}{% comments
}

\content{
        \begin{example}
        You can afford a \$700 per month mortgage payment. You’ve found a 30 year loan at
        5\% interest.
        \begin{enumerate}
                \item How big of a loan can you afford?
                \item How much total money will you pay the loan company?
                \item How much of that money is interest?        
        \end{enumerate}
        \end{example}
}{% presentation
\vspace{3in}
}{% nopresentation
\vspace{3in}
}{% comments
}

\content{
        \begin{example}
        You deposit \$1000 each year into an account earning 8\% compounded annually.
        \begin{enumerate}
                \item How much will you have in the account in 10 years?
                \item How much total money will you put into the account?
                \item How much total interest will you earn?        
        \end{enumerate}
        \end{example}
}{% presentation
}{% nopresentation
}{% comments
}
